\section{Datasets, Labels, and Leakage Controls}
\label{sec:data}

\subsection{Internal Development Data}
The internal v0.2 corpus supports training, validation, and development-time evaluation. Its examples are assembled from deterministic templates and public-source-derived security cases. The data mixture deliberately includes hard negatives---requests that contain suspicious tokens but should remain safe---and an explicit unknown class for cases in which model-visible context is insufficient. Prompt-injection examples target both the verdict and the recommended-action fields.

Internal evaluation is useful for checking training convergence and output alignment, but it is not the primary evidence in this paper. On the sealed 1,750-example independent split, the LoRA model achieved 0.9657 verdict accuracy, 1.000 malicious recall, 0.0286 normal FPR, and 1.000 unknown accuracy. It nevertheless failed three locked release gates: JSON syntax success was 0.9714, schema validity was 0.9714, and prompt-injection obedience was 0.9197. Eleven of 137 prompt-injection cases hijacked the verdict or action recommendation. These results motivated an external challenge evaluation rather than a production-readiness claim.

\subsection{Frozen \challenge{} Challenge Set}
The primary test set contains 413 authorized LAN request-side samples. Its model-input file has SHA-256
\begin{center}
{\small\texttt{c77f4949\allowbreak cbe1860f\allowbreak 97219a83\allowbreak f0ee16de\allowbreak 40740321\allowbreak a4d2f628\allowbreak d56f1506\allowbreak 9fe6b8c3}}.
\end{center}
The frozen reference-label distribution is 209 safe, 97 attack, and 107 unknown. The label file has SHA-256 \labelhash. The set is external to the v0.2 train, validation, closed-test, and independent-test splits. It is not used for training, prompt selection, threshold selection, hyperparameter tuning, checkpoint selection, or model selection.

The challenge set contains request-side fields only. Labels and scoring metadata are stored separately. Cohorts include normal requests, attack attempts, hard negatives, prompt-injection cases, and insufficient-context examples. The dataset is an authorized challenge capture, not a random sample of natural production traffic; prevalence-sensitive metrics must therefore not be interpreted as production base rates.

\subsection{AI-Assisted Annotation and Adjudication}
Two independent AI annotators produced initial labels. Agreement was high but imperfect, with observed agreement 0.9274 and Cohen's $\kappa=0.8651$. A third AI model adjudicated disagreements. A separate model-visible-context audit corrected 68 provenance issues in which labels depended on information that the classifier could not observe. The resulting frozen reference set contains 160 direct-consensus outcomes and 253 model-adjudicated outcomes. The reference file was frozen before the AetherArgus challenge predictions were generated.

The AI-assisted stage is not described as human ground truth. It is a reproducible label-construction process with known dependence on model judgment. To reduce this threat, the full frozen reference set was subsequently subjected to a confidential blind human review.

\subsection{Confidential Blind Human Review}
Human review occurred after the frozen model evaluation. Reviewers operated outside the public repository in a restricted environment and made their decisions before seeing annotator-A or annotator-B labels, the third-model adjudication, AetherArgus predictions, or evaluation metrics. After blind decisions were complete, the human outcomes were compared with the frozen references.

The review covered all 413 samples, found zero disagreements, and required zero label changes. The before- and after-review label hashes are therefore identical. No prediction rerun or rescoring was required. The repository records \texttt{humanReviewed=true}, full coverage, zero disagreements, and a SHA-bound operator attestation.

Reviewer identities, reviewer count, worksheets, and raw materials are withheld because they contain confidential benchmark and operational information. This limits external auditability. We consequently use the term \emph{confidentially human-reviewed reference labels}, not publicly reproducible human ground truth.

\subsection{Frozen Protocol and Process Isolation}
Protocol v1 was committed before the rule baseline. Protocol v2 was committed before the remaining seven primary model runs and fixes model identities, data hashes, metrics, statistical tests, system-prompt hash, few-shot example hash, model-asset hashes, and deterministic generation parameters.

Inference and scoring are process-separated:
\begin{equation}
    \text{model inputs} \rightarrow \text{predictions.jsonl}
\end{equation}
\begin{equation}
    \text{predictions.jsonl} + \text{labels} \rightarrow \text{metrics}.
\end{equation}
The inference programs do not open the label file or adjudication artifacts. Predictions are append-only, flushed per sample, and keyed by unique sample identifiers. The strict validator checks 413-of-413 coverage, duplicate and omission absence, common input hash, prediction-label alignment, train/test separation, few-shot exclusion, metric recomputation, figure/table consistency, privacy scans, and claim-to-evidence traceability. The frozen evidence package passed all 27 strict checks.
